%%%%%%%%%%%%%%%%%%%%%%%%%%%%%%%%%%%%%%%%%%%%%%%%%%%%%%%%%%%%
% UNCOMMENT THE LINE BELOW TO VIEW COMMENTS
% \def\DraftVersion {}
%%%%%%%%%%%%%%%%%%%%%%%%%%%%%%%%%%%%%%%%%%%%%%%%%%%%%%%%%%%%

% TITLE TO BE PRINTED ON THE HEADERS
\newcommand{\titleOnHeader}{\textsf{IMITATOR} \emph{not} a Developer Manual}

% TITLE TO BE PRINTED ON THE FIRST PAGE
\newcommand{\titleOnFirstPage}{IMITATOR \emph{not} a Developer Manual}

\input{commons.tex}

\newcommand{\validator}{\textsf{VALIDATOR}}

%%%%%%%%%%%%%%%%%%%%%%%%%%%%%%%%%%%%%%%%%%%%%%%%%%%%%%%%%%%%
%%%%%%%%%%%%%%%%%%%%%%%%%%%%%%%%%%%%%%%%%%%%%%%%%%%%%%%%%%%%
%%%%%%%%%%%%%%%%%%%%%%%%%%%%%%%%%%%%%%%%%%%%%%%%%%%%%%%%%%%%

% META DATA
\hypersetup{
	pdftitle={IMITATOR Developer Manual},
	pdfauthor={\'Etienne Andr\'e}%
}
\title{IMITATOR Developer Manual}
\author{Étienne André}
	


%%%%%%%%%%%%%%%%%%%%%%%%%%%%%%%%%%%%%%%%%%%%%%%%%%%%%%%%%%%%
%%%%%%%%%%%%%%%%%%%%%%%%%%%%%%%%%%%%%%%%%%%%%%%%%%%%%%%%%%%%
%%%%%%%%%%%%%%%%%%%%%%%%%%%%%%%%%%%%%%%%%%%%%%%%%%%%%%%%%%%%
\begin{document}
% \maketitle


%%%%%%%%%%%%%%%%%%%%%%%%%%%%%%%%%%%%%%%%%%%%%%%%%%%%%%%%%%%%
\input{title_and_toc.tex}
%%%%%%%%%%%%%%%%%%%%%%%%%%%%%%%%%%%%%%%%%%%%%%%%%%%%%%%%%%%%




%%%%%%%%%%%%%%%%%%%%%%%%%%%%%%%%%%%%%%%%%%%%%%%%%%%%%%%%%%%%
%%%%%%%%%%%%%%%%%%%%%%%%%%%%%%%%%%%%%%%%%%%%%%%%%%%%%%%%%%%%
\chapter{Introduction}
%%%%%%%%%%%%%%%%%%%%%%%%%%%%%%%%%%%%%%%%%%%%%%%%%%%%%%%%%%%%
%%%%%%%%%%%%%%%%%%%%%%%%%%%%%%%%%%%%%%%%%%%%%%%%%%%%%%%%%%%%

\imitator{}~\cite{Andre21} has no developer manual, for many (mainly bad) reasons:
lack of time, small team of developers with strong interactions, little concurrent development.

Nevertheless, this document tries (or will try) to summarize some of the important internal characteristics of \imitator{}.
It will naturally evolve over time.


%%%%%%%%%%%%%%%%%%%%%%%%%%%%%%%%%%%%%%%%%%%%%%%%%%%%%%%%%%%%
%%%%%%%%%%%%%%%%%%%%%%%%%%%%%%%%%%%%%%%%%%%%%%%%%%%%%%%%%%%%
\chapter{Development environment}
%%%%%%%%%%%%%%%%%%%%%%%%%%%%%%%%%%%%%%%%%%%%%%%%%%%%%%%%%%%%
%%%%%%%%%%%%%%%%%%%%%%%%%%%%%%%%%%%%%%%%%%%%%%%%%%%%%%%%%%%%

The OCaml version used to develop \imitator{} is OCaml 4.08.0.


%%%%%%%%%%%%%%%%%%%%%%%%%%%%%%%%%%%%%%%%%%%%%%%%%%%%%%%%%%%%
%%%%%%%%%%%%%%%%%%%%%%%%%%%%%%%%%%%%%%%%%%%%%%%%%%%%%%%%%%%%
\chapter{Parsing process}
%%%%%%%%%%%%%%%%%%%%%%%%%%%%%%%%%%%%%%%%%%%%%%%%%%%%%%%%%%%%
%%%%%%%%%%%%%%%%%%%%%%%%%%%%%%%%%%%%%%%%%%%%%%%%%%%%%%%%%%%%


Since version 3.0, the parsing process is as follows:

\begin{enumerate}
	\item Get the options from the terminal command
	\item Parse the model
	\item Depending on the option (notably, the presence or not of a \styleOption{-mode}, that triggers translation or syntax check instead of looking for a property), decide whether a property should be parsed or not; if so, parse the property
	\item Check and convert to abstract structures at the same time both the model and the property (if any)
\end{enumerate}



%%%%%%%%%%%%%%%%%%%%%%%%%%%%%%%%%%%%%%%%%%%%%%%%%%%%%%%%%%%%
%%%%%%%%%%%%%%%%%%%%%%%%%%%%%%%%%%%%%%%%%%%%%%%%%%%%%%%%%%%%
\chapter{Constraint representation}
%%%%%%%%%%%%%%%%%%%%%%%%%%%%%%%%%%%%%%%%%%%%%%%%%%%%%%%%%%%%
%%%%%%%%%%%%%%%%%%%%%%%%%%%%%%%%%%%%%%%%%%%%%%%%%%%%%%%%%%%%

Constraints are implemented using the Parma Polyhedra Library (PPL)~\cite{BHZ08}.
Constraints in PPL use integer-valued indexes for variables.

In \imitator{}, the variables are numbered as follows in the constraint, starting from~0:
\begin{itemize}
	\item first the parameters
	\item then the regular clocks in the model
	\item then the observer clock (if any), automatically inferred from the property (if any)
	\item then the extra reset clock (mainly used to measure time elapsing in non-Zeno checking)
	\item then all discrete variables.
\end{itemize}

% TODO: dessin avec les variables et une accolade dessous montrant les indices




%%%%%%%%%%%%%%%%%%%%%%%%%%%%%%%%%%%%%%%%%%%%%%%%%%%%%%%%%%%%
%%%%%%%%%%%%%%%%%%%%%%%%%%%%%%%%%%%%%%%%%%%%%%%%%%%%%%%%%%%%
\chapter{Type inference}
%%%%%%%%%%%%%%%%%%%%%%%%%%%%%%%%%%%%%%%%%%%%%%%%%%%%%%%%%%%%
%%%%%%%%%%%%%%%%%%%%%%%%%%%%%%%%%%%%%%%%%%%%%%%%%%%%%%%%%%%%

Since version 3.1 and the addition of \styleIMI{int} and \styleIMI{bool} types, a type inference system is added to \imitator{}.
Details can be obtained on how the types are inferred using the following command:

%------------------------------------------------------------
\begin{terminal}
imitator model.imi -verbose high | grep "Infer"
\end{terminal}
%------------------------------------------------------------


%%%%%%%%%%%%%%%%%%%%%%%%%%%%%%%%%%%%%%%%%%%%%%%%%%%%%%%%%%%%
%%%%%%%%%%%%%%%%%%%%%%%%%%%%%%%%%%%%%%%%%%%%%%%%%%%%%%%%%%%%
\chapter{Algorithms class diagrams}
%%%%%%%%%%%%%%%%%%%%%%%%%%%%%%%%%%%%%%%%%%%%%%%%%%%%%%%%%%%%
%%%%%%%%%%%%%%%%%%%%%%%%%%%%%%%%%%%%%%%%%%%%%%%%%%%%%%%%%%%%

We give below the class diagram of the non-distributed algorithms implemented in \imitator{} below.
Virtual classes are depicted in yellow.


\includegraphics[width=\textwidth]{diagrams/classDiagramSimplified.png}



The full class diagram of all algorithm classes (included distributed algorithms) is given in the next page.

\newenvironment{changemargin}[2]{%
\begin{list}{}{%
\setlength{\topsep}{0pt}%
\setlength{\leftmargin}{#1}%
\setlength{\rightmargin}{#2}%
\setlength{\listparindent}{\parindent}%
\setlength{\itemindent}{\parindent}%
\setlength{\parsep}{\parskip}%
}%
\item[]}{\end{list}}

\newpage
\begin{landscape}
% \newgeometry{left=3cm,bottom=0.1cm}
\begin{changemargin}{-2.cm}{-3cm}
\thispagestyle{empty}
\includegraphics[width=2.3\textwidth]{diagrams/classDiagramFull.png}
\end{changemargin}
\end{landscape}
\newpage



%%%%%%%%%%%%%%%%%%%%%%%%%%%%%%%%%%%%%%%%%%%%%%%%%%%%%%%%%%%%%
%%%%%%%%%%%%%%%%%%%%%%%%%%%%%%%%%%%%%%%%%%%%%%%%%%%%%%%%%%%%%
\chapter{Testing}
%%%%%%%%%%%%%%%%%%%%%%%%%%%%%%%%%%%%%%%%%%%%%%%%%%%%%%%%%%%%%
%%%%%%%%%%%%%%%%%%%%%%%%%%%%%%%%%%%%%%%%%%%%%%%%%%%%%%%%%%%%%

Non-regression tests can be called using


%------------------------------------------------------------
\begin{terminal}
python tests/test.py
\end{terminal}
%------------------------------------------------------------

It is possible to call only a subset of the tests using the tag system.
An example of command is:

%------------------------------------------------------------
\begin{terminal}
python tests/test.py --filter="tags=mydesiredtag"
\end{terminal}
%------------------------------------------------------------

\noindent{}where ``\texttt{mydesiredtag}'' represents the tag the tests should match to be called.

The non-regression tests are defined in the file \texttt{tests/regression\_tests\_data.py}.


%%%%%%%%%%%%%%%%%%%%%%%%%%%%%%%%%%%%%%%%%%%%%%%%%%%%%%%%%%%%
\section{Random testing with \validator{}}
%%%%%%%%%%%%%%%%%%%%%%%%%%%%%%%%%%%%%%%%%%%%%%%%%%%%%%%%%%%%

\validator{} is a tool for developers that looks for bugs in \imitator{} using random models.
It generates random parametric timed automata, runs \imitator{} on each of them with several configurations (without starting a new process), and reports a model on which the results differ.
The model is then reduced to a small one that still shows the difference, and saved as an \texttt{.imi} file.
\validator{} is not part of the \imitator{} releases; its sources are in the folder \texttt{validator/}.

%------------------------------------------------------------
\subsection{Building}
%------------------------------------------------------------

\validator{} needs the library \texttt{qcheck-core}, which is not needed to build \imitator{}; for this reason it is not built by \texttt{dune build}.

%------------------------------------------------------------
\begin{terminal}
opam install qcheck-core
dune build @validator
\end{terminal}
%------------------------------------------------------------

\noindent{}The binary is \texttt{bin/validator}.

%------------------------------------------------------------
\subsection{Running}
%------------------------------------------------------------

\validator{} takes a specification file (see below) and some options:

%------------------------------------------------------------
\begin{terminal}
./bin/validator validator/example.val -r 200 -time-limit 2
\end{terminal}
%------------------------------------------------------------

\begin{description}
	\item[\texttt{-output DIR}] folder in which the models are saved (default: \texttt{validator-output}).
	\item[\texttt{-r N}] maximum number of models to generate and test (default: 1000).
	\item[\texttt{-time-limit S}] time limit in seconds for each run of \imitator{} (decimals are allowed). A model on which a run reaches the time limit is skipped; so is a model on which a run is stopped by a depth or states limit given in a configuration.
	\item[\texttt{-seed N}] seed of the random generator. The same seed and specification give the same models, which makes it possible to reproduce a run. Without this option, a seed is chosen at random and printed.
	\item[\texttt{-results-file FILE}] write a summary of the search in JSON.
\end{description}

Paths in the specification file are relative to the folder from which \validator{} is called: the example has to be run from the root of the repository.

%------------------------------------------------------------
\subsection{Specification file}
%------------------------------------------------------------

A specification file has two sections: \texttt{[mode]} says what to do, and \texttt{[model]} describes the models to generate.
Each line is of the form \texttt{key = value}; lines starting with \texttt{\#} are comments.
For example (see also \texttt{validator/example.val}):

\begin{verbatim}
[mode]
kind = compare
configs = [
    tests/testcases/basic-properties/synth-parametric-win-strat-accepting.imiprop
    -PTG-no-strategy-generation;

    tests/testcases/basic-properties/synth-parametric-win-strat-accepting.imiprop
    -PTG-no-strategy-generation
    -PTG-abstraction ch
]

[model]
nb_automata = 1
nb_clocks = 1..2
nb_parameters = 1..2
nb_locations = 3..8
transitions_per_location = 1..2
controllability_ratio = 0.5
\end{verbatim}

\paragraph{Section \texttt{[model]}}
A number is either an integer, or a range \texttt{lo..hi} (both bounds included) in which a value is chosen for each model.
The keys \texttt{nb\_locations}, \texttt{nb\_clocks} and \texttt{nb\_parameters} are mandatory, except in the mode \texttt{reduce}.
There must be at least one automaton, location, clock and transition per location.

\begin{description}
	\item[\texttt{nb\_automata}] number of automata (default: 1).
	\item[\texttt{nb\_locations}] number of locations of each automaton.
	\item[\texttt{nb\_clocks}, \texttt{nb\_parameters}] number of clocks and of parameters.
	\item[\texttt{transitions\_per\_location}] number of outgoing transitions of each location (default: 1).
	\item[\texttt{all\_reachable}] if \texttt{true}, every location is reachable from the initial location in the graph of the automaton (default: \texttt{true}).
	\item[\texttt{cycles}] if \texttt{false}, the graph of the automaton has no cycle (default: \texttt{true}).
	\item[\texttt{guard\_probability}, \texttt{invariant\_probability}] probability that a transition has a guard, and that a location has an invariant (default: 0.5).
	\item[\texttt{reset\_probability}] probability that a transition resets clocks (default: 0.2).
	\item[\texttt{controllability\_ratio}] probability that a transition is controllable, for parametric timed games (default: 0).
	\item[\texttt{max\_constant}] largest constant used in guards and invariants (default: 10).
\end{description}

The automata are named \texttt{automaton\_0}, \texttt{automaton\_1}, etc., and their locations \texttt{l0}, \texttt{l1}, etc.; \texttt{l0} is the initial location.
The clocks are named \texttt{x}, \texttt{y}, \texttt{z}, \texttt{clk\_3}, etc., and the parameters \texttt{p}, \texttt{q}, \texttt{r}, \texttt{p\_3}, etc.
These names can be used in the property files of the configurations.
A property must only refer to locations that exist in all generated models: with \texttt{nb\_locations = 3..8}, a property can refer to \texttt{l0}, \texttt{l1} and \texttt{l2}.

\paragraph{Section \texttt{[mode]}}
The key \texttt{kind} selects one of the following modes (default: \texttt{sample}).

\begin{description}
	\item[\texttt{compare}]
		Run \imitator{} on each model once per configuration, and stop at the first model on which the resulting constraints differ.
		The key \texttt{configs} is a list of at least two configurations between brackets, separated by semicolons.
		A configuration is the command line of \imitator{} without the model: a property file followed by options.
		The model and its reduced version are saved in the subfolder \texttt{counter\_examples}, as \texttt{counter\_example\_<id>.imi} and \texttt{counter\_example\_<id>\_reduced.imi}.

	\item[\texttt{reduce}]
		Reduce an existing model on which the configurations of \texttt{configs} give different results.
		The key \texttt{model} is the path of the model; the section \texttt{[model]} is not needed.
		The reduced model is saved in the subfolder \texttt{reducer}, as \texttt{<model>\_reduced.imi}.

	\item[\texttt{sample}]
		Only generate models, to look at what a \texttt{[model]} section produces.
		The key \texttt{samples} is the number of models (default: 1); with \texttt{pdf = true}, each model is also drawn using \texttt{dot}.
		The models are saved in the subfolder \texttt{samples}, as \texttt{sampled\_model\_1.imi}, etc.
\end{description}

%------------------------------------------------------------
\subsection{Reading the results}
%------------------------------------------------------------

In the mode \texttt{compare}, \validator{} prints the number of the current model, and the number of models skipped because of a time limit (\texttt{TO}).
When a model is found, it is saved, then reduced: transitions are removed, guards and invariants are replaced with \texttt{True}, and locations are merged, as long as the difference between the configurations remains.
The reduced model is saved next to the full one, with the suffix \texttt{\_reduced}.
Running \imitator{} on the reduced model with each configuration shows the difference.
If no model is found after the number of models given by \texttt{-r}, nothing is saved: this does not mean that the configurations are equivalent.

A difference is not always a bug: two configurations may legitimately give different results (for example two different properties, or an option that makes the result an over-approximation).
\validator{} compares the constraints of the results only, not their soundness.

If \imitator{} raises an exception, returns an error or aborts on a generated model, or if the results cannot be compared, \validator{} stops and saves this model in the subfolder \texttt{crashes}, as \texttt{crash\_<id>.imi}.
During the reduction, a model on which this happens is only discarded.

The exit code is 0 when the search ends, whether a model is found or not.
It is 1 in case of an error: an invalid specification file or configuration, a failure on a model, or nothing to reduce in the mode \texttt{reduce}.

The comparison of the results is defined in \texttt{validator/comparison/Comparison.ml}; only the results made of a single constraint, or of a good and a bad constraint, are supported so far.

%------------------------------------------------------------
\subsection{Maintaining \validator{}}
%------------------------------------------------------------

\validator{} is not built by the continuous integration: a change in \imitator{} can break its compilation without anyone noticing.
It was last tested with \texttt{qcheck-core} 0.91.

The sources are organized as follows:

\begin{description}
	\item[\texttt{VALIDATOR.ml}] the entry point, with one case per mode. The other files at the root of the folder are the command line (\texttt{ValidatorOptions}), the loading of the configurations (\texttt{ConfigLoader}), the search (\texttt{CounterExampleFinder}), the saving of models (\texttt{ModelOutput}) and the statistics on the generated guards (\texttt{CoverageAnalyzer}).
	\item[\texttt{config/}] the parser of the specification files.
	\item[\texttt{gen/}] the generators: graphs of the automata (\texttt{Adjacency}), guards and invariants (\texttt{PZone}, \texttt{PZoneGen}), models (\texttt{SimpleModelGen}), and their translation into the parsing structures of \imitator{} (\texttt{Convert}).
	\item[\texttt{runner/}] the run of \imitator{} on a generated model (\texttt{ModelRunner}).
	\item[\texttt{comparison/}] the comparison of two results.
	\item[\texttt{reducer/}] the reduction; \texttt{ModelReducer} applies the steps one after the other.
	\item[\texttt{printer/}] the messages printed in the terminal.
\end{description}

\texttt{ModelRunner} converts a generated model and runs the algorithm selected by \texttt{AlgorithmDriver}, as \texttt{IMITATOR.ml} does, but many times in the same process.
\imitator{} keeps some global data between two runs, which \texttt{ModelRunner} has to reset: the dimensions of the constraints (\texttt{PVal}, \texttt{HyperRectangle}, \texttt{LinearConstraint}) and the cache of invariants (\texttt{State}).
If a run seems to depend on the runs before it, look for such data first.

To compare a new type of result, add a case to \texttt{eq\_result} in \texttt{Comparison.ml}.
To add a mode, add a constructor to \texttt{config/Mode.ml}, parse it in \texttt{config/Parser.ml}, and add a case to \texttt{VALIDATOR.ml}.

To check the generator itself after changing it, there is one more mode, \texttt{kind = coverage}: it generates guards only (their number is given by \texttt{-r}, and their shape by \texttt{nb\_clocks}, \texttt{nb\_parameters} and \texttt{max\_constant}) and prints statistics on the kinds of constraints obtained, to see whether they are varied enough.
With \texttt{-results-file}, the statistics are also written in JSON.
The mode \texttt{sample} is the other way to look at what is generated.


%%%%%%%%%%%%%%%%%%%%%%%%%%%%%%%%%%%%%%%%%%%%%%%%%%%%%%%%%%%%%
%%%%%%%%%%%%%%%%%%%%%%%%%%%%%%%%%%%%%%%%%%%%%%%%%%%%%%%%%%%%%
\chapter{Licensing and Credits}
%%%%%%%%%%%%%%%%%%%%%%%%%%%%%%%%%%%%%%%%%%%%%%%%%%%%%%%%%%%%%
%%%%%%%%%%%%%%%%%%%%%%%%%%%%%%%%%%%%%%%%%%%%%%%%%%%%%%%%%%%%%


\subsection*{Developer Manual}
This developer manual is available under the Creative Commons CC-BY-SA license.

\begin{center}
	\includegraphics[width=.2\textwidth]{CC-BY-SA_500.png}
\end{center}


%%%%%%%%%%%%%%%%%%%%%%%%%%%%%%%%%%%%%%%%%%%%%%%%%%%%%%%%%%%%%
\input{graphics-credits.tex}

\input{bibliography.tex}
%%%%%%%%%%%%%%%%%%%%%%%%%%%%%%%%%%%%%%%%%%%%%%%%%%%%%%%%%%%%%




\end{document}

